\documentclass[11pt,a4paper]{article}

\usepackage[T1]{fontenc}
\usepackage[utf8]{inputenc}
\usepackage[spanish]{babel}
\usepackage[a4paper,margin=2.8cm]{geometry}
\usepackage{microtype}
\usepackage{xurl}
\usepackage[hidelinks]{hyperref}

\hypersetup{
  pdftitle={Todo ángel es terrible},
  pdfauthor={Juan José Gómez Cadenas}
}

\setlength{\parindent}{1.25em}
\setlength{\parskip}{0.35em}
\raggedbottom

\title{Todo ángel es terrible}
\author{Juan José Gómez Cadenas}
\date{}

\begin{document}

\maketitle

En su brillante artículo, «Las tijeras de Queneau en la Biblioteca de Babel», ejecuta mi estimadísimo Ángel Fernández un curioso ejercicio estilístico, copiando a quien rebate. Todavía está caliente en la prensa virtual de Jot Down mi reciente artículo sobre Borges en su biblioteca, cuando la rauda pluma de mi amigo produce una respuesta que lo replica, sacándose de la manga un espejo del país de las maravillas. La Biblioteca de Babel encuentra su imagen en la máquina de Queneau y la condena a la que el pobre Borges se ve abocado ---leer basura sin sentido por toda la eternidad--- es alterada por un aguerrido franchute que irrumpe en la biblioteca pidiendo unas tijeras. Como hacía mi yo juvenil con Rilke, como dice el protagonista sin nombre de «Hombre de la esquina rosada», el artículo de Ángel copia al mío hasta en la forma de escupir. Y por supuesto el plagio es intencionado y parte del juego.

Sostiene Ángel que la creación no está en inventar el texto sino en seleccionarlo. Cito literalmente:

\begin{quotation}
«Le roi de la pampa retourne sa chemise / depuis que lord Elgin négligea ses naseaux»

«El rey de la pampa cambia de camisa / desde que lord Elgin descuidó sus ollares» no parece el comienzo de un poema de Neruda. La combinación de personajes parece arbitraria en la forma en que trastoca y construye las palabras y, aunque es clara como una escritura notarial, no parece llevarnos a la exaltación del espíritu que uno espera de la poesía.
\end{quotation}

Anota, precisa lectora, sofisticado lector, este detalle: \emph{no parece llevarnos a la exaltación del espíritu que uno espera de la poesía.} Volveré más tarde sobre ello.

El misterio de la críptica frase se explica enseguida:

\begin{quotation}
[El libro en cuestión es] \emph{Cent mille milliards de poèmes}, publicado en 1961 por el escritor Raymond Queneau. Este señor francés no solo era matemático, sino también uno de los fundadores del Oulipo ---un taller de escritura cuyos miembros creían que se podía mejorar la calidad de la literatura imponiéndole reglas que cualquier persona sensata consideraría motivo justificado para dejar el oficio---.

El libro recoge diez sonetos, uno por página, y cada página está dividida en catorce tiras, en cada una de las cuales solo hay un verso, para que puedas unir el primer verso de un soneto con el segundo de otro ---como en esos libros infantiles en los que un bombero acaba con cabeza de jirafa---, con la promesa de que, sea cual sea la combinación, se mantendrán el ritmo, la rima y la métrica correcta. Así, en el primer soneto, el rey de la pampa cuelga su camisa de los cuernos de un toro para que se seque; y en el segundo, el caballo del Partenón ---el verdadero dueño de los ollares--- se indigna por la negligencia de lord Elgin; de este modo, basta con cambiar una sola tira para convertir el saqueo de la Acrópolis de Atenas en un asunto de lavandería.
\end{quotation}

Así, el juego de Queneau propone combinar las tiras de versos a capricho del lector, ofreciendo un repertorio casi inagotable ---la friolera de cien billones--- de posibles sonetos. Frente al horror combinatorio de la Biblioteca, Ángel y Raimundo nos ofrecen ilimitadas posibilidades creativas... simplemente combinando texto generado por otro(s). Si reemplazamos al autor de las tiras por una IA, tenemos una justificación perfectamente válida de coescritura entre hombre y máquina.

Es un espléndido argumento que, sin embargo, no contradice en absoluto nada de lo que yo he escrito antes.

\emph{Primo}, nada tengo contra la coautoría hombre-máquina. Solo reclamo que se le dé a César lo que es de César y a Claude lo que es de Claude. El problema de los escritores pillados con las manos en la masa ---o debería decir en la IA--- en recientes concursos literarios es que han hecho trampas, saltándose las reglas. Me parece excelente, como proponen algunos escritores, que se convoquen concursos híbridos y se publiquen novelas en las que Claude sea coautor(a). Les auguro, eso sí, el mismo éxito que han tenido los torneos híbridos de ajedrez. En ese bello juego-deporte que tanto me apasiona ---y en el que tan torpe soy---, se siguen con devoción las partidas entre humanos y se hace poco o ningún caso a las competiciones entre máquinas o entre equipos hombre-máquina. Extrañamente, a pesar de que las máquinas juegan mejor que nosotros, no nos interesan. El ajedrez no es solo un juego. Es un combate de boxeo psicológico, un drama silencioso, un teatro de kabuki, una batalla mortal en la que los contendientes no se levantan de la silla, una exhibición de belleza, un catálogo de errores, una explosión de intuición, un torrente de emociones reprimidas. AlphaZero puede jugar todo lo bien que quiera, pero no nos emociona. Dudo mucho que las novelas o los poemas híbridos (o puramente IA) que tanto defienden mis amigos funcionen mejor. AlphaZero juega mejor que un humano porque el ajedrez se puede reducir a un algoritmo programable en una máquina. Pero la poesía no es un algoritmo y quien crea que un LLM puede capturar las \emph{Elegías}, o no ha leído a Rilke, o no ha leído las mediocridades que producen los chatbots.

\emph{Secondo}, el texto de Ángel pasa Pangram, con un 99~\% humano, a pesar de que su estructura de espejo y su estilo podrían inducir sospechas de que el autor ha usado la IA para «copiar» con maestría el artículo con el que se bate ---una sospecha que podría agudizarse si tomamos en cuenta que ha escrito el artículo en menos de 48 horas---. La provocación no puede estar más clara. Afirma Ángel que \emph{cualquiera con suficiente oficio puede tomarle el pelo a Pangram}. Y predica con el ejemplo.

Excepto que eso no está tampoco en contradicción con nada que yo haya escrito. Como ya he repetido muchas veces ---y esta es la última que vuelvo sobre el tema---, Pangram está optimizado para concedernos el beneficio de la duda y su gran utilidad es detectar texto obviamente IA, no descubrir exquisitas construcciones de exquisitos falsificadores. La clave es detener el torrente de \emph{slop} que producen innumerables teclados incontinentes, y los artesanos capaces de marear a Pangram son tan escasos ---o más escasos si cabe--- como los escritores que se niegan a delegar en una máquina.

La comparación con el ajedrez viene otra vez al caso. No hay cosa más absurda que jugar una partida online ayudándose de un motor. La gracia del ajedrez es, precisamente, jugar. Cuando un tramposo deja que el robot mueva por él las piezas, le está haciendo perder el tiempo a su oponente y se está engañando a sí mismo en la misma medida que un halterófilo que hiciera sentadillas con 500 kilos... ayudado por una máquina que se los levantara. Y aun así, el ajedrez online está infectado de tramposos. En los clubs online, por supuesto, se pasan con regularidad programas para detectarlos, que, sin embargo, no cazan a todo el mundo. Hay jugadores ---sobre todo entre los que echan partidas largas o por turnos, de ahí que nadie quiera hacerlo ya--- que tienen la paciencia de equivocarse a propósito de vez en cuando, elegir sistemáticamente la segunda o tercera mejor jugada que sugiere la máquina en lugar de la primera, todo para estafar al detector y ganar la partida... ¿con qué propósito? Y sin embargo, los clubs online siguen pasando los detectores \emph{para desanimar a la mayoría de los tramposos}, esto es, para mitigar el \emph{slop}. Mi argumento para favorecer Pangram no es que piense que se puede detectar a todos los tramposos ---y repito, llamo tramposo al que escribe con IA sin declararlo--- sino en desanimar a los ejércitos de narcisos armados de un chatbot.

\emph{Terzo} y más sustancial. Todo el complejo y espejado artículo con el que Ángel nos deleita parece estar discutiendo con un \emph{eidolon}, una imagen mía conjurada por alguna divinidad con ganas de broma, no conmigo. El ejercicio de Queneau y compañía me parece divertidísimo en la categoría de juegos florales, interesante conceptualmente ---en el mismo sentido que el intento de Hilbert de formalizar las matemáticas a principios del siglo pasado, convirtiéndolas en una máquina formal de producir teoremas---. He escrito en otro artículo:

\begin{quotation}
No afirmo que un texto generado o cogenerado por IA carezca de interés. Tengo muchos ejemplos de lo contrario. Sostengo, sin embargo, firme, apasionadamente, que las lágrimas de las cosas no están en ellos.
\end{quotation}

El propio Ángel, en el arranque de su artículo, se pega un tiro en el pie:

\begin{quotation}
y, aunque es clara como una escritura notarial, no parece llevarnos a la exaltación del espíritu que uno espera de la poesía.
\end{quotation}

Exactamente. Y ese es mi argumento. Las combinaciones que produce Queneau, o las que produce un LLM bien entrenado, o las que produce un equipo hombre-máquina, merecen justamente el epíteto notarial, pero no se cuentan en todo lo que merece lágrimas.

Merecen lágrimas la muerte de Dido y la de Julieta, la escena en la que la Maga descubre que su bebé ha dejado de respirar, el momento en que un personaje de \emph{Cien años de soledad} atraviesa un páramo verde donde el eco repite los pensamientos, Príamo besando las manos de Aquiles, el momento en el que Orfeo mira hacia atrás y pierde para siempre a Eurídice, los versos de Dylan Thomas, «Rage, rage against the dying of the light», la tenacidad del escribiente que decide no hacerlo, la locura de Ahab, el amor desesperado del padre que lleva una sola bala en su revólver para matar a su hijo antes que abandonarlo a las alimañas que pueblan la Tierra.

Sostengo que ninguna IA ---no las que conocemos--- escribirá líneas comparables a esas. Y la razón la explica muy bien Ángel. Una cosa es divertirse reordenando renglones escritos en tiras de papel y otra cosa es captar la extraña cualidad del alma humana. Que el alma existe es evidente: la literatura no es otra cosa que la comunión entre ellas, las almas de los vivos y de los muertos. Que el alma se aniquila con el final de la vida también lo es para mí, aunque mis amigos cristianos opinen otra cosa. Que esa tragedia ---\emph{rage, rage against the dying of the light}--- es la esencia central de la literatura, no me cabe duda. Somos seres de carne y hueso, de piel que ha sentido otra piel rozándonos, de ojos que se han maravillado contemplando el océano, de brazos que han sostenido a nuestros bebés y manos que han acariciado a nuestros amantes, somos sacos de protoplasma con un kilo y medio de materia gris en el cráneo y los intestinos llenos de mierda, y es la suma de esa mierda y esos sesos donde se conectan un billón de conexiones y esa piel que tanto recuerda y esa boca que ha besado y mordido y ese deseo que nos enloquece y nos eleva, la que nos convierte en seres espirituales, ángeles cargados de cadenas, capaces de entender, \emph{de intuir} la belleza de lo efímero y lo transcendente, como ningún LLM ---armado solo de palabras--- y ningún ejército de monos locos o pedantes armados con tijeras podría jamás entender:

\begin{verse}
Pues la belleza no es otra cosa que el principio del terror,\\
y nos sorprende cuando, serenamente,\\
desdeña aniquilarnos.\\
Todo ángel es terrible.
\end{verse}

\end{document}
